\documentclass[11pt,a4paper]{article}
\usepackage[utf8]{inputenc}
\usepackage[top=0.9in, bottom=0.8in, left=1in, right=1in]{geometry}
\usepackage{hyperref}

\setlength{\parskip}{0.8em}
\setlength{\parindent}{0pt}

\begin{document}

\begin{center}
    {\LARGE \textbf{Liuxi Mei}} \\
    \vspace{2mm}
    Linköping, Sweden $\cdot$ liume102@student.liu.se $\cdot$ +46-0704898962
    \vspace{4mm}
\end{center}

\textbf{Research Statement} \\
\textbf{Position:} Doctoral Student in Machine Learning \\
\textbf{Project:} Machine learning for sensing in distributed wireless systems \\
\textbf{Date:} April 9, 2026

\section*{Research Motivation}

I am applying for this project because it connects machine learning with physical hardware. A D-MIMO system generates large sets of structured data. The antennas and radio links form a specific physical layout. A practical model should utilize this geometry rather than just mapping inputs to outputs. It must also process the data fast enough for real-time sensing.

My master's thesis led me to this approach. I used graph neural networks (GNNs) to study biological networks. Biology and wireless communication have different applications, but the underlying mathematics is similar. They are both systems of connected parts. A D-MIMO system can be modeled as a physical graph. The antennas are the nodes, and the radio links are the edges. GNNs are well-suited for this setup. Instead of sending all the raw data to a central server, each antenna can communicate with its neighbors. This local message passing reduces latency and network traffic. It is a practical way to make the sensing system fast and scalable.

\section*{Research Goals}

Deeply understanding the hardware sensing problem is the first major step. The wireless data changes frequently. Finding the exact reasons for these changes is critical. Testing basic models will show where they fail. This process will reveal when catching an anomaly becomes difficult.

Building models that use the network structure is another key focus. These new models need to work well in unseen environments. Making efficient use of the distributed nodes is essential. The system must also remain robust when the data gets noisy.

Testing these models against real hardware limits is equally important. High accuracy on a clean dataset is insufficient. The real cost in network speed and computational power is a major constraint. Proving that these methods stay stable on an actual testbed is the primary goal.

\section*{Suggested Thesis Contents}

The contents of my thesis will cover several core topics around machine learning for distributed sensing. The work is organized into the following four areas.

\textbf{Studying Testbed Data.} To begin with, the data from the D-MIMO testbed needs close attention. The radio signals change over time. The antennas connect to each other in specific ways. The environment also changes the data. These patterns must be understood before selecting a model. This phase will identify the main challenges of learning in this setup.

\textbf{Applying Flexible Graph Models.} After that, the focus will shift to exploring GNNs. Graph models do not depend on a fixed number of nodes. New antennas can be added or removed, and the trained model can still maintain performance. This flexibility is highly beneficial for real hardware setups. My thesis provided experience with graph attention models to guide the flow of information smoothly. This will help the system detect specific events more efficiently.

\textbf{Building Context-Aware Methods.} Context-aware models will also be designed. In my master's work, complex cell contexts were modeled successfully. A similar approach will be applied to D-MIMO sensing. The model will learn the stable background as one part, and the sudden changes as another part. This separation helps the model find true anomalies in new environments. It will not require complete retraining every time the weather or the room changes.

\textbf{Optimizing Local Computation.} Fast and local computing is the final step. Working at Ericsson provided a practical understanding of network congestion. Sending all raw data to a central computer causes significant delays. Models will be designed to use local message passing, where antennas communicate only with their close neighbors. This approach will keep the sensing accurate while significantly reducing the network load.

\end{document}